% ================================================================
% LỜI GIẢI ĐỀ THI 1: GIẢI TÍCH HÀM NHIỀU BIẾN
% Lời giải chi tiết với các bước giải rõ ràng
% ================================================================

\section*{Lời giải Đề 1: Giải tích hàm nhiều biến}
\addcontentsline{toc}{section}{Lời giải Đề 1: Giải tích hàm nhiều biến}

% ================================================================
% CÂU 1A: ĐẠO HÀM RIÊNG VÀ ĐỊNH LÝ EULER
% ================================================================
\subsection*{Câu 1a: Đạo hàm riêng và Định lý Euler}

\begin{questionbox}
Cho hàm số $u(x, y) = (x + y)\sqrt{x^2 + y^2}$. Chứng minh rằng: 
\[
    x\dfrac{\partial u}{\partial x} + y\dfrac{\partial u}{\partial y} = 2u 
    \quad \text{với mọi } (x, y) \neq (0,0).
\]
\end{questionbox}

\begin{solution}
Cho hàm số: $u(x, y) = (x + y)\sqrt{x^2 + y^2}$

Cần chứng minh hệ thức: $x\dfrac{\partial u}{\partial x} + y\dfrac{\partial u}{\partial y} = 2u$

\textit{} Tính đạo hàm riêng theo $x$:
\begin{align*}
    \dfrac{\partial u}{\partial x} &= \dfrac{\partial}{\partial x}(x+y) \cdot \sqrt{x^2 + y^2} + (x + y) \cdot \dfrac{\partial}{\partial x}(\sqrt{x^2 + y^2}) \\
    &= 1 \cdot \sqrt{x^2 + y^2} + (x + y) \cdot \dfrac{x}{\sqrt{x^2 + y^2}} \\
    &= \dfrac{(x^2 + y^2) + x(x + y)}{\sqrt{x^2 + y^2}} \\
    &= \dfrac{2x^2 + y^2 + xy}{\sqrt{x^2 + y^2}}
\end{align*}

\textit{} Tính đạo hàm riêng theo $y$:
\begin{align*}
    \dfrac{\partial u}{\partial y} &= \dfrac{\partial}{\partial y}(x+y) \cdot \sqrt{x^2 + y^2} + (x + y) \cdot \dfrac{\partial}{\partial y}(\sqrt{x^2 + y^2}) \\
    &= 1 \cdot \sqrt{x^2 + y^2} + (x + y) \cdot \dfrac{y}{\sqrt{x^2 + y^2}} \\
    &= \dfrac{(x^2 + y^2) + y(x + y)}{\sqrt{x^2 + y^2}} \\
    &= \dfrac{x^2 + 2y^2 + xy}{\sqrt{x^2 + y^2}}
\end{align*}

\textit{} Kiểm tra hệ thức:
\begin{align*}
    x\dfrac{\partial u}{\partial x} + y\dfrac{\partial u}{\partial y} &= \dfrac{x(2x^2 + y^2 + xy) + y(x^2 + 2y^2 + xy)}{\sqrt{x^2 + y^2}} \\
    &= \dfrac{2x^3 + xy^2 + x^2y + x^2y + 2y^3 + xy^2}{\sqrt{x^2 + y^2}} \\
    &= \dfrac{2x^3 + 2xy^2 + 2x^2y + 2y^3}{\sqrt{x^2 + y^2}} \\
    &= \dfrac{2(x+y)(x^2+y^2)}{\sqrt{x^2 + y^2}} \\
    &= 2(x + y)\sqrt{x^2 + y^2} = 2u
\end{align*}

\textit{Kết luận:} Hệ thức đã được chứng minh. $\blacksquare$
\end{solution}

% ================================================================
% CÂU 1B: TÌM CỰC TRỊ TỰ DO
% ================================================================
\subsection*{Câu 1b: Tìm cực trị tự do}

\begin{questionbox}
Tìm cực trị của hàm số: $z = (x + y)^2 - \dfrac{x^4}{2} - \dfrac{y^4}{2}$.
\end{questionbox}

\begin{solution}
Cho hàm số: $z = (x + y)^2 - \dfrac{x^4}{2} - \dfrac{y^4}{2}$

\textit{} Tìm các điểm dừng bằng cách giải hệ phương trình đạo hàm riêng bậc nhất:
\begin{align*}
    z'_x &= 2(x + y) - 2x^3 = 0 \\
    z'_y &= 2(x + y) - 2y^3 = 0
\end{align*}

Từ hệ trên suy ra: $2x^3 = 2y^3 \Rightarrow x^3 = y^3 \Rightarrow x = y$.

Thay vào phương trình đầu: $2(2x) - 2x^3 = 0 \Rightarrow 4x - 2x^3 = 0 \Rightarrow 2x(2 - x^2) = 0$

Suy ra: $x = 0$ hoặc $x^2 = 2 \Rightarrow x = \pm\sqrt{2}$.

\textit{Các điểm dừng:} $M_1(0, 0)$, $M_2(\sqrt{2}, \sqrt{2})$, $M_3(-\sqrt{2}, -\sqrt{2})$.

\textit{} Kiểm tra điều kiện đủ bằng ma trận Hessian:
\begin{align*}
    A = z''_{xx} &= 2 - 6x^2 \\
    C = z''_{yy} &= 2 - 6y^2 \\
    B = z''_{xy} &= 2
\end{align*}

Định thức Hessian: $D = AC - B^2$

\textit{Tại $M_1(0, 0)$:}
$A = 2, C = 2, B = 2 \Rightarrow D = 2 \cdot 2 - 2^2 = 0$.

Vì $D = 0$, cần kiểm tra thêm:
- Dọc theo $y = x$: $z(x, x) = 4x^2 - x^4$, có cực tiểu tại $x = 0$.
- Dọc theo $y = -x$: $z(x, -x) = -x^4$, có cực đại tại $x = 0$.

Vậy $M_1(0, 0)$ là \textit{điểm yên ngựa}.

\textit{Tại $M_2(\sqrt{2}, \sqrt{2})$ và $M_3(-\sqrt{2}, -\sqrt{2})$:}
$A = 2 - 6(2) = -10, C = -10, B = 2$
$D = (-10)(-10) - 2^2 = 100 - 4 = 96 > 0$.

Vì $A < 0$ và $D > 0$, hàm số đạt \textit{cực đại địa phương} tại hai điểm này.

Giá trị cực đại: 
\begin{align*}
    z(M_2) = z(M_3) &= (2\sqrt{2})^2 - \dfrac{(\sqrt{2})^4}{2} - \dfrac{(\sqrt{2})^4}{2} \\
    &= 8 - \dfrac{4}{2} - \dfrac{4}{2} = 8 - 2 - 2 = 4
\end{align*}

\textit{Kết luận:} 
- Điểm yên ngựa: $(0, 0)$ với $z = 0$
- Cực đại địa phương: $(\pm\sqrt{2}, \pm\sqrt{2})$ với $z = 4$
\end{solution}

% ================================================================
% CÂU 1C: CỰC TRỊ CÓ ĐIỀU KIỆN
% ================================================================
\subsection*{Câu 1c: Cực trị có điều kiện}

\begin{questionbox}
Tìm cực trị của hàm số $z = 20x^{1/4}y^{3/4}$, thỏa mãn điều kiện $x + 3y = 8$ với $x, y > 0$.
\end{questionbox}

\begin{solution}
Cho hàm số: $z = 20x^{1/4}y^{3/4}$ với điều kiện $x + 3y = 8$ ($x, y > 0$).

\textit{Phương pháp nhân tử Lagrange:}
Xét hàm Lagrange: $L(x, y, \lambda) = 20x^{1/4}y^{3/4} - \lambda(x + 3y - 8)$

\textit{} Tính các đạo hàm riêng:
\begin{align*}
    L'_x &= 5x^{-3/4}y^{3/4} - \lambda = 0 \\
    L'_y &= 15x^{1/4}y^{-1/4} - 3\lambda = 0 \\
    L'_\lambda &= -(x + 3y - 8) = 0
\end{align*}

\textit{} Từ hai phương trình đầu:
\begin{align*}
    \lambda &= 5x^{-3/4}y^{3/4} = 5\left(\dfrac{y}{x}\right)^{3/4} \\
    \lambda &= 5x^{1/4}y^{-1/4} = 5\left(\dfrac{x}{y}\right)^{1/4}
\end{align*}

Cân bằng hai biểu thức của $\lambda$:
\[
    5\left(\dfrac{y}{x}\right)^{3/4} = 5\left(\dfrac{x}{y}\right)^{1/4}
\]

\[
    \left(\dfrac{y}{x}\right)^{3/4} = \left(\dfrac{x}{y}\right)^{1/4} = \left(\dfrac{y}{x}\right)^{-1/4}
\]

\[
    \left(\dfrac{y}{x}\right)^{3/4 + 1/4} = 1 \Rightarrow \dfrac{y}{x} = 1 \Rightarrow y = x
\]

\textit{} Thay vào điều kiện ràng buộc:
$x + 3x = 8 \Rightarrow 4x = 8 \Rightarrow x = 2$

Vậy $x = 2, y = 2$.

\textit{} Tính giá trị cực trị:
\[
    z(2, 2) = 20(2)^{1/4}(2)^{3/4} = 20 \cdot 2^{1/4 + 3/4} = 20 \cdot 2 = 40
\]

\textit{Kết luận:} Giá trị lớn nhất là $40$ tại điểm $(2, 2)$.
\end{solution}

% ================================================================
% CÂU 2A: TÍCH PHÂN KÉP BẰNG CÁCH ĐỔI BIẾN
% ================================================================
\subsection*{Câu 2a: Tích phân kép bằng cách đổi biến}

\begin{questionbox}
Tính diện tích hình phẳng $D$ giới hạn bởi các đường:
\[
    y = \dfrac{1}{x}, \quad y = \dfrac{5}{x}, \quad x = y^2, \quad x = 5y^2.
\]
\end{questionbox}

\begin{solution}
Miền $D$ được giới hạn bởi: $xy = 1$, $xy = 5$, $\dfrac{x}{y^2} = 1$, $\dfrac{x}{y^2} = 5$.

\textit{} Đặt phép đổi biến:
\[
    u = xy, \quad v = \dfrac{x}{y^2}
\]

Khi đó miền $D$ trở thành: $D' = [1, 5] \times [1, 5]$ trong hệ tọa độ $(u, v)$.

\textit{} Tìm biểu thức của $x, y$ theo $u, v$:
Từ $u = xy$ và $v = \dfrac{x}{y^2}$, ta có:
\[
    \dfrac{u}{v} = \dfrac{xy}{\dfrac{x}{y^2}} = xy \cdot \dfrac{y^2}{x} = y^3
\]

Suy ra: $y = \left(\dfrac{u}{v}\right)^{1/3} = u^{1/3}v^{-1/3}$

Và: $x = \dfrac{u}{y} = \dfrac{u}{u^{1/3}v^{-1/3}} = u^{2/3}v^{1/3}$

\textit{} Tính định thức Jacobian:
\[
    J = \dfrac{\partial(x,y)}{\partial(u,v)} = \begin{vmatrix}
        \dfrac{\partial x}{\partial u} & \dfrac{\partial x}{\partial v} \\
        \dfrac{\partial y}{\partial u} & \dfrac{\partial y}{\partial v}
    \end{vmatrix}
\]

\begin{align*}
    \dfrac{\partial x}{\partial u} &= \dfrac{2}{3}u^{-1/3}v^{1/3}, \quad 
    \dfrac{\partial x}{\partial v} = \dfrac{1}{3}u^{2/3}v^{-2/3} \\
    \dfrac{\partial y}{\partial u} &= \dfrac{1}{3}u^{-2/3}v^{-1/3}, \quad 
    \dfrac{\partial y}{\partial v} = -\dfrac{1}{3}u^{1/3}v^{-4/3}
\end{align*}

\begin{align*}
    |J| &= \left| \dfrac{2}{3}u^{-1/3}v^{1/3} \cdot \left(-\dfrac{1}{3}u^{1/3}v^{-4/3}\right) - \dfrac{1}{3}u^{2/3}v^{-2/3} \cdot \dfrac{1}{3}u^{-2/3}v^{-1/3} \right| \\
    &= \left| -\dfrac{2}{9}v^{-1} - \dfrac{1}{9}v^{-1} \right| = \left| -\dfrac{3}{9v} \right| = \dfrac{1}{3v}
\end{align*}

\textit{} Tính diện tích:
\begin{align*}
    S &= \iint_D dx\,dy = \iint_{D'} |J| \, du\,dv \\
    &= \int_1^5 \int_1^5 \dfrac{1}{3v} \, du\,dv \\
    &= \dfrac{1}{3} \int_1^5 du \int_1^5 \dfrac{1}{v} \, dv \\
    &= \dfrac{1}{3} \cdot 4 \cdot \ln 5 = \dfrac{4\ln 5}{3}
\end{align*}

\textit{Kết luận:} Diện tích hình phẳng $D$ là $\dfrac{4\ln 5}{3}$.
\end{solution}

% ================================================================
% CÂU 2B: BÀI TOÁN MẬT ĐỘ DÂN SỐ
% ================================================================
\subsection*{Câu 2b: Bài toán mật độ dân số}

\begin{questionbox}
Mật độ dân số của một thành phố được cho bởi hàm $f(x, y) = 20000e^{-(x^2+y^2)/5}$, tính theo mặt phẳng tọa độ $Oxy$ với gốc là trung tâm thành phố (đơn vị: người/km²). Tính số dân của thành phố trong khu vực có bán kính 5 km tính từ trung tâm thành phố. Mật độ dân số trung bình của khu vực trên là bao nhiêu?
\end{questionbox}

\begin{solution}
Mật độ dân số: $f(x, y) = 20000e^{-(x^2+y^2)/5}$ (người/km²)

\textit{} Tính tổng dân số trong hình tròn $x^2 + y^2 \leq 25$:
\[
    P = \iint_{x^2+y^2 \leq 25} 20000e^{-(x^2+y^2)/5} \, dA
\]

\textit{} Sử dụng tọa độ cực:
$x = r\cos\theta$, $y = r\sin\theta$, $dA = r \, dr \, d\theta$

Với $0 \leq r \leq 5$ và $0 \leq \theta \leq 2\pi$:
\[
    P = \int_0^{2\pi} d\theta \int_0^5 20000e^{-r^2/5} \cdot r \, dr
\]

\textit{} Tính tích phân theo $r$:
Đặt $t = \dfrac{r^2}{5} \Rightarrow dt = \dfrac{2r}{5} dr \Rightarrow r \, dr = \dfrac{5}{2} dt$

Khi $r = 0 \Rightarrow t = 0$; khi $r = 5 \Rightarrow t = 5$:
\begin{align*}
    \int_0^5 20000e^{-r^2/5} \cdot r \, dr &= 20000 \cdot \dfrac{5}{2} \int_0^5 e^{-t} \, dt \\
    &= 50000 \left[ -e^{-t} \right]_0^5 \\
    &= 50000(-e^{-5} + 1) \\
    &= 50000(1 - e^{-5})
\end{align*}

\textit{} Tính tổng dân số:
\begin{align*}
    P &= 2\pi \cdot 50000(1 - e^{-5}) \\
    &= 100000\pi(1 - e^{-5})
\end{align*}

Với $e^{-5} \approx 0.00674$:
\[
    P \approx 100000\pi(1 - 0.00674) \approx 100000\pi \cdot 0.99326 \approx 311,895 \text{ người}
\]

\textit{} Tính mật độ dân số trung bình:
Diện tích khu vực: $S = \pi \cdot 5^2 = 25\pi$ km²

Mật độ trung bình:
\begin{align*}
    \bar{f} &= \dfrac{P}{S} = \dfrac{100000\pi(1 - e^{-5})}{25\pi} \\
    &= 4000(1 - e^{-5}) \\
    &\approx 4000 \times 0.99326 \approx 3,973 \text{ người/km}^2
\end{align*}

\textit{Kết luận:}
- Tổng dân số: $100000\pi(1 - e^{-5}) \approx 311,895$ người
- Mật độ trung bình: $4000(1 - e^{-5}) \approx 3,973$ người/km²
\end{solution}

% ================================================================
% CÂU 3: TÍNH KHỐI LƯỢNG VẬT THỂ
% ================================================================
\subsection*{Câu 3: Tính khối lượng vật thể}

\begin{questionbox}
Trong không gian $\mathbb{R}^3$, cho vật thể:
\[
    \Omega = \left\{ (x, y, z) \in \mathbb{R}^3 : \sqrt{x^2 + y^2} \leq z \leq 4 - \sqrt{x^2 + y^2} \right\}
\]
có khối lượng riêng phân bố theo hàm $\rho(x, y, z) = x^2 + y^2$. Tính khối lượng của $\Omega$.
\end{questionbox}

\begin{solution}
Vật thể $\Omega$ giới hạn bởi hai mặt nón $z = r$ và $z = 4 - r$ với $r = \sqrt{x^2+y^2}$.

\textit{} Tìm miền xác định:
Giao tuyến của hai mặt nón: $r = 4 - r \Rightarrow 2r = 4 \Rightarrow r = 2$.

Vậy $0 \leq r \leq 2$ và $r \leq z \leq 4-r$.

\textit{} Sử dụng tọa độ trụ:
$x = r\cos\theta$, $y = r\sin\theta$, $z = z$

Với $0 \leq \theta \leq 2\pi$, $0 \leq r \leq 2$, $r \leq z \leq 4-r$.

Khối lượng riêng: $\rho = x^2 + y^2 = r^2$

Phần tử thể tích: $dV = r \, dr \, d\theta \, dz$

\textit{} Tính khối lượng:
\begin{align*}
    m &= \iiint_\Omega \rho \, dV \\
    &= \int_0^{2\pi} d\theta \int_0^2 r \, dr \int_r^{4-r} r^2 \, dz \\
    &= \int_0^{2\pi} d\theta \int_0^2 r^3 (4-r-r) \, dr \\
    &= 2\pi \int_0^2 r^3 (4 - 2r) \, dr \\
    &= 2\pi \int_0^2 (4r^3 - 2r^4) \, dr \\
    &= 2\pi \left[ r^4 - \dfrac{2r^5}{5} \right]_0^2 \\
    &= 2\pi \left( 16 - \dfrac{2 \cdot 32}{5} \right) \\
    &= 2\pi \left( 16 - \dfrac{64}{5} \right) \\
    &= 2\pi \cdot \dfrac{80 - 64}{5} \\
    &= 2\pi \cdot \dfrac{16}{5} = \dfrac{32\pi}{5}
\end{align*}

\textit{Kết luận:} Khối lượng của vật thể $\Omega$ là $\dfrac{32\pi}{5}$.
\end{solution}

% ================================================================
% CÂU 4: TÍCH PHÂN ĐƯỜNG VÀ ĐỊNH LÝ GREEN
% ================================================================
\subsection*{Câu 4: Tích phân đường và Định lý Green}

\begin{questionbox}
Tính tích phân đường:
\[
    I = \int_\gamma (5y - 6^x)\,dx + (10x + 6^y)\,dy,
\]
trong đó $\gamma$ là đường tròn $x^2 + y^2 = 36$.
\end{questionbox}

\begin{solution}
Tính tích phân đường: $I = \oint_\gamma (5y - 6^x) \, dx + (10x + 6^y) \, dy$

với $\gamma$ là đường tròn $x^2 + y^2 = 36$ định hướng ngược chiều kim đồng hồ.

\textit{} Áp dụng định lý Green:
Với $P(x,y) = 5y - 6^x$ và $Q(x,y) = 10x + 6^y$:
\[
    I = \iint_D \left( \dfrac{\partial Q}{\partial x} - \dfrac{\partial P}{\partial y} \right) dA
\]

\textit{} Tính các đạo hàm riêng:
\begin{align*}
    \dfrac{\partial Q}{\partial x} &= \dfrac{\partial}{\partial x}(10x + 6^y) = 10 \\
    \dfrac{\partial P}{\partial y} &= \dfrac{\partial}{\partial y}(5y - 6^x) = 5
\end{align*}

\textit{} Áp dụng công thức Green:
\begin{align*}
    I &= \iint_D (10 - 5) \, dA = 5 \iint_D dA \\
    &= 5 \times \text{Diện tích}(D)
\end{align*}

\textit{} Tính diện tích miền $D$:
Miền $D$ là hình tròn có bán kính $R = 6$ (vì $x^2 + y^2 = 36$).

Diện tích: $\text{Area}(D) = \pi R^2 = \pi \cdot 6^2 = 36\pi$

\textit{} Kết quả cuối cùng:
\[
    I = 5 \times 36\pi = 180\pi
\]

\textit{Kết luận:} Tích phân đường $I = 180\pi$.
\end{solution}

\newpage